\documentclass[11pt]{article}
\usepackage[margin=1in]{geometry}
\usepackage{amsmath,amssymb}
\usepackage{bm}
\usepackage[T1]{fontenc}
\usepackage{hyperref}
\usepackage{enumitem}

\title{FiberForm Porosity Codebase Overview}
\author{Codebase Notes}
\date{\today}

\begin{document}
\maketitle

\section{Purpose}
This note gives a current overview of the code in
\texttt{porosity\_code}. The present codebase supports three main covariance
workflow stages, along with a few supporting porosity utilities:
\begin{enumerate}[label=\arabic*.]
\item empirical covariance construction from voxel data,
\item covariance-model fitting,
\item continuous porosity-field generation on a simulation domain.
\end{enumerate}

The document also records important caveats that matter when interpreting the
outputs.

\section{High-Level Workflow}
At a high level, the code takes a binarized TIFF stack and can:
\begin{enumerate}[label=\arabic*.]
\item read the binary volume,
\item compute global porosity and a blockwise porosity field,
\item compute empirical covariance of the porosity field,
\item fit an isotropic exponential covariance model,
\item generate continuous Gaussian porosity realizations from the fitted model.
\end{enumerate}

The entry point is \texttt{main.c}. Input parsing is handled by
\texttt{input.c} and \texttt{input.h}.

\section{Input Structure}
Each stack is currently supplied through an \texttt{INPUT\_DATA} entry of the
form:
\begin{verbatim}
INPUT_DATA = "/path/to/stack.tif 2.0" ;
FILTER_RESOLUTION = 50.0
\end{verbatim}

Here:
\begin{itemize}
\item the path points to a TIFF stack,
\item the number in the same quoted entry is that stack's voxel resolution,
\item \texttt{FILTER\_RESOLUTION} is shared across the run.
\end{itemize}

Multiple stacks are supported by giving multiple quoted entries in
\texttt{INPUT\_DATA}.

\section{Run Modes}
The code now supports staged execution through explicit switches:
\begin{verbatim}
RUN_FULL_PIPELINE = YES|NO
EMPIRICAL_COVARIANCE = YES|NO
FIT_COVARIANCE_MODEL = YES|NO
GENERATE_POROSITY_REALIZATIONS = YES|NO
\end{verbatim}

Additional utility switches still exist:
\begin{itemize}
\item \texttt{GLOBAL\_POROSITY}
\item \texttt{POROSITY\_FIELD}
\item \texttt{POROSITY\_COVARIANCE\_LAG}
\item \texttt{POROSITY\_AXIS}
\item \texttt{POROSITY\_COVARIANCE\_REF\_WINDOW}
\item \texttt{REF\_WINDOW\_VALUES}
\item \texttt{POROSITY\_REALIZATION\_COUNT}
\item \texttt{POROSITY\_REALIZATION\_GRID\_RES}
\item \texttt{POROSITY\_REALIZATION\_GRID\_DIM}
\end{itemize}

\section{Stage 0: TIFF Reading and Binary Interpretation}
The TIFF reading helpers are implemented in \texttt{tiff.c}. This code:
\begin{itemize}
\item opens 8-bit single-channel TIFF pages,
\item verifies that all slices in the stack have consistent dimensions,
\item loads the full volume into memory when needed,
\item supports both direct grayscale indices and palette TIFFs.
\end{itemize}

The binary conversion helper is \texttt{palette\_index\_to\_binary(...)}.

\subsection{Important Current Convention}
The active code path currently treats binary value \(0\) as void in the actual
porosity counting and porosity-field construction logic.

This means the live behavior is effectively:
\[
0 = \text{void}, \qquad \text{nonzero} = \text{solid or non-void}.
\]

This convention must be kept in mind when interpreting porosity numbers.

\section{Global Porosity}
Global porosity is handled by \texttt{porosity\_global.c}. The code:
\begin{enumerate}[label=\arabic*.]
\item reads the full stack statistics,
\item counts total voxels,
\item counts void voxels,
\item reports porosity as
\[
\phi_{\mathrm{global}} = \frac{N_{\mathrm{void}}}{N_{\mathrm{total}}}.
\]
\end{enumerate}

If more than one stack is processed in the same run, the program also prints the
average porosity across stacks.

\section{Local Porosity Field}
Local porosity-field construction is implemented in \texttt{porosity\_field.c}.

\subsection{What It Computes}
The code converts the binary volume into a continuous-valued porosity field by
coarse-graining the volume into blocks whose size is determined by:
\[
\texttt{FILTER\_RESOLUTION} / \texttt{VOXEL\_RESOLUTION}.
\]

The filter size in voxels is rounded to an odd integer. Then the code:
\begin{enumerate}[label=\arabic*.]
\item builds a 3D integral volume for fast block summation,
\item partitions the domain into non-overlapping blocks,
\item computes the void fraction in each block,
\item writes the resulting porosity field as a floating-point TIFF stack.
\end{enumerate}

\subsection{Interpretation}
This is a blockwise coarse-grained porosity field, not a sliding-window field.
That distinction matters later because the covariance is the covariance of this
coarse-grained field.

\section{Stage 1: Empirical Covariance}
Empirical covariance construction is implemented in
\texttt{porosity\_covariance\_lag.c}.

\subsection{Mathematical Form}
Once the porosity field \(\phi(\bm{x})\) is available, the code computes its
mean \(\mu\) and variance \(\sigma^2\), and then estimates covariance at lag
\(\bm{r}\) by a finite-sample average over valid overlapping pairs:
\[
C(\bm{r}) =
\frac{1}{N(\bm{r})}
\sum_{\bm{x}\in\Omega_{\bm{r}}}
(\phi(\bm{x})-\mu)(\phi(\bm{x}+\bm{r})-\mu).
\]

The implementation then normalizes that quantity by the variance:
\[
\rho(\bm{r}) = \frac{C(\bm{r})}{\sigma^2}.
\]

So the stored covariance table is effectively a normalized covariance or
correlation-like quantity.

\subsection{Supported Modes}
The empirical covariance stage supports:
\begin{itemize}
\item \texttt{3DLAG}: the full nonnegative 3D lag grid,
\item \texttt{CURVE}: a covariance curve along a selected axis.
\end{itemize}

\subsection{Outputs}
Stage 1 writes outputs into a sibling folder named:
\begin{verbatim}
BUILD_COVARIANCE
\end{verbatim}

Typical outputs are:
\begin{itemize}
\item a text lag table,
\item a VTK file for ParaView.
\end{itemize}

\section{Stage 2: Covariance Preparation and Fitting}
This stage is implemented in \texttt{porosity\_covariance\_realization\_prep.c}
and described by the public types in
\texttt{porosity\_covariance\_realization\_prep.h}.

\subsection{What It Does}
This stage performs more than a bare fit. It:
\begin{enumerate}[label=\arabic*.]
\item reads the \texttt{3DLAG} table,
\item reconstructs the raw covariance from the normalized covariance,
\item symmetrizes the covariance grid,
\item analyzes tail noise,
\item analyzes negative covariance lobes,
\item analyzes anisotropy,
\item fits an isotropic exponential covariance model,
\item validates the fit,
\item evaluates the fitted covariance on the requested simulation grid,
\item validates the evaluated simulation-grid covariance.
\end{enumerate}

\subsection{Model Form}
The current fit is isotropic exponential:
\[
C(r) = s \exp\left(-\frac{r}{\ell}\right),
\]
where \(s\) is the sill and \(\ell\) is the correlation length.

\subsection{Outputs}
Stage 2 also writes into \texttt{BUILD\_COVARIANCE}. It produces:
\begin{itemize}
\item a tail-noise report,
\item an exponential-fit quality report,
\item a covariance-build summary,
\item a fitted-model artifact file.
\end{itemize}

The fitted-model artifact is the handoff into Stage 3.

\section{Stage 3: Continuous Porosity Realization Generation}
Generation is implemented in
\texttt{porosity\_covariance\_realization\_generation.c}.

\subsection{Method}
The current generator:
\begin{enumerate}[label=\arabic*.]
\item reads the fitted exponential covariance model,
\item builds the covariance matrix on the requested simulation grid,
\item performs Cholesky factorization,
\item draws Gaussian random vectors,
\item generates continuous porosity realizations,
\item writes both TIFF and VTK outputs.
\end{enumerate}

If the fitted model is \(C(\bm{r})\), then the generated field is a Gaussian
field whose covariance is prescribed by that model and whose mean is taken from
the fitted-model artifact.

\subsection{Outputs}
Generated realizations are written into:
\begin{verbatim}
GENERATED_POROSITY_FIELDS
\end{verbatim}

Each realization is written in:
\begin{itemize}
\item floating-point TIFF format,
\item VTK format for ParaView.
\end{itemize}

\section{Reference-Window Covariance Utility}
There is also an auxiliary module,
\texttt{porosity\_covariance\_refwindow.c}, which supports a
reference-window-based covariance analysis. This is not the main route used by
the current three-stage realization workflow, but it remains available as an
additional diagnostic path.

\section{Current Caveats}
The following caveats are important:
\begin{itemize}
\item The current \texttt{README.md} had previously been stale relative to the
live code and needed updating.
\item The active porosity convention in the live code path currently behaves as
\(0 = \text{void}\).
\item The porosity field is block-based, not sliding-window.
\item The fit stage is broader than pure model fitting because it also checks
simulation-grid compatibility.
\item The realization generator uses a full covariance matrix and Cholesky
factorization, so it is not scalable to large grids.
\item The covariance model currently used for generation is isotropic
exponential only.
\end{itemize}

\section{Summary}
In its present form, the codebase is a staged porosity-to-covariance-to-field
pipeline. It begins with binarized TIFF stacks, computes a coarse-grained local
porosity field, estimates empirical covariance of that field, fits an isotropic
exponential covariance model, and finally generates continuous Gaussian porosity
realizations on a user-defined simulation grid.

\end{document}
